\chapter{Hamilton--Jacobi 可达性分析}
\label{ch:hj-reachability}

% ════════════════════════════════════════════════════════════════
%  P4 控制部 · C1b 第 I 章 · 归卷四《控制理论基础》
%  主源：Tutorial 01_数学/40_控制理论/160_HJ可达性分析.md（2076 行）
%  设计稿：docs/控制部设计_C1b_I.md（§4 大纲 / §2 勘误 / §3 记号 / §5 cref / §6 bib）
%  头号去重边界：与 P3 ch:plan-diffgame（已写）—— HJI/黏性解/BRS-BRT/Air3D 一律 \cref 不重抄；
%  本章 = diff-game 提名的四增量：CBVF / DeepReach 深推导 / 数值双曲 PDE 理论 / 解耦框架。
%  CBVF 主证在本章 §I.5（F 章 clf_cbf_safety.tex eq:cbvf 给桥、引本章主证）。
% ════════════════════════════════════════════════════════════════

\noindent
本章把“集合级、最坏情况、无限时域的安全验证”讲透：从二人零和微分博弈出发的 Hamilton--Jacobi--Isaacs（HJI）变分不等式，给出从动力学\emph{自动算出}的最优安全集；再沿四条主线深化——后向可达管的障碍算子、水平集数值方法的双曲 PDE 理论、把控制屏障函数与 HJ 可达性折扣统一的控制屏障-值函数（CBVF），以及用神经网络突破维度灾难的 DeepReach。它在卷四《控制理论基础》中承上启下：为 \cref{ch:clf-cbf} 的控制屏障函数提供一个“算出来的、最优的”屏障函数，为 \cref{ch:dp-hjb} 的 HJB 黏性解理论给出对抗推广，并把 \cref{ch:lyapunov-stability} 的吸引域、\cref{ch:robust-hinf} 的最坏情况博弈统一进同一台“值函数 + 水平集”机器。
\medskip

% ----------------------------------------------------------------
\section*{承上：从“会收敛”“永远安全”“鲁棒”到“可证安全集”}

卷四已沿三条线建立了控制证书。\cref{ch:lyapunov-stability} 用 Lyapunov 函数证明系统\emph{会收敛}；\cref{ch:clf-cbf} 用控制屏障函数证明系统\emph{永不离开}容许集；\cref{ch:robust-hinf} 用 $H_\infty$/博弈视角处理\emph{最坏情况}干扰。三者都需要先手中握有一个证书函数——Lyapunov 函数、屏障函数 $h$、或鲁棒增益。但证书从何而来？低维可手工凑，高维、强非线性、带对抗干扰的系统则无从下手。本章给出一条根本不同的路：不去\emph{猜}证书，而是把“从哪些状态出发、在最坏干扰下不可避免地出事”这个集合\emph{算}出来——这个集合的边界本身就是最严格、最不保守的安全证书，而它对应的值函数就是一个现成的、最优的控制屏障函数。代价是维度灾难，这也是后半章数值方法与神经逼近着力的地方。

\subsection*{本章与微分博弈章（\texorpdfstring{\cref{ch:plan-diffgame}}{差分博弈章}）的分工：四个增量}

HJ 可达性的奠基理论——零和博弈到 HJI 方程的推导、Hamiltonian、Isaacs 条件、黏性解、后向可达集（BRS）/后向可达管（BRT）/reach--avoid 的定义、Air3D 与 Homicidal Chauffeur 经典追逃——已在规划部 \cref{ch:plan-diffgame} 完整建立。\textbf{凡这些定义与主推导，本章一律交叉引用，绝不重写}：HJI 方程见 \cref{eq:dg-hji}、Hamiltonian 见 \cref{eq:dg-ham}、Isaacs 条件见 \cref{eq:dg-isaacs-cond,eq:dg-separable}、值函数即唯一黏性解见 \cref{thm:dg-evans} 与 \cref{def:dg-viscosity}、BRS/BRT/reach--avoid 见 \cref{eq:dg-brs-def,eq:dg-brs-pde,eq:dg-brt-def,eq:dg-brt-vi,eq:dg-brt-mitchell,eq:dg-reachavoid}、Lax--Friedrichs 数值 Hamiltonian 见 \cref{eq:dg-lf}、Air3D 相对动力学见 \cref{eq:dg-air3d}、维度诅咒与通往 CBF 的接口见 \cref{sec:dg-curse,sec:dg-frs-cbf}。本章只做 \cref{ch:plan-diffgame} 明确“留待控制部”的四块深化：

\begin{enumerate}
  \item \textbf{BRT 障碍算子的动态规划起源}（\cref{sec:hj-brt-obstacle}）——把 $\min\{0,H\}$ 那个看似神秘的截断项从动态规划原理一步推出来，并补一个 Isaacs 不可分离的具体反例；
  \item \textbf{水平集数值方法的双曲 PDE 理论}（\cref{sec:hj-numerics,sec:hj-cfl}）——迎风差分“为什么必须”、LF 隐式迎风等价于数值扩散、CFL 条件的严格推导、维度灾难的量化；规划部偏“写一个求解器”的工程口径，本章补偏微分方程理论视角；
  \item \textbf{CBVF：CBF 与 HJ 的折扣统一}（\cref{sec:hj-cbvf}）——本章第一理论高峰，\cref{ch:plan-diffgame} 完全没有，直接服务 \cref{ch:clf-cbf}；
  \item \textbf{DeepReach 与解耦式安全框架}（\cref{sec:hj-deepreach,sec:hj-frameworks}）——高维神经逼近的深推导，与 FaSTrack/Fisac/多智能体分解的完整数学。
\end{enumerate}

\begin{note}
\textbf{本章记号约定（与 \cref{ch:plan-diffgame} 对齐）.}\;
状态 $x\in\mathbb R^n$，值函数 $V(x,t)$，时间偏导记 $D_tV:=\partial V/\partial t$（与 \cref{eq:dg-brt-mitchell} 一致），空间梯度 $\nabla_xV$，协态 $\lambda=\nabla_xV$（影子价格，见 \cref{sec:dg-bridge}）。
\textbf{控制 $u\in\mathcal U$、扰动 $d\in\mathcal D$}——本章沿用控制部主流记号；与 \cref{ch:plan-diffgame} 可达节文献记号的对应是 $a\leftrightarrow u$（控制）、$b\leftrightarrow d$（扰动）。
目标/危险集用 Lipschitz 隐式函数编码，$\mathcal T=\{x:g(x)\le0\}$，$g$ 为带符号距离（即 \cref{eq:dg-sdf} 的 $g$，源文献常记 $l$，本书统一为 $g$ 以与 \cref{ch:plan-diffgame} 一致）。安全集取零超水平集 $\{V\ge0\}$（与 CBVF 原文一致）。Hamiltonian $H(x,\lambda)=\max_{u}\min_{d}\lambda\cdot f(x,u,d)$ 即 \cref{eq:dg-ham}。
\textbf{两处同字母消歧}：$\alpha_i$ 仅指 Lax--Friedrichs 的第 $i$ 维数值黏性系数（限 \cref{sec:hj-numerics,sec:hj-cfl}），与扩展类 $\mathcal K_\infty$ 函数 $\alpha(\cdot)$（CBF 衰减，限 \cref{sec:hj-cbvf}、定义见 \cref{def:ext-kinf}）作用域隔离；$\omega_0$ 仅指 SIREN 的频率超参（限 \cref{sec:hj-deepreach}），与角速度无关。CBVF 值函数记 $B_\gamma$，折扣参数 $\gamma\ge0$。
\end{note}

% ════════════════════════════════════════════════════════════════
\section{为什么需要可达性：屏障函数的边界与 HJ 的承诺}
\label{sec:hj-why}

\paragraph{动机：手工屏障函数的三重困难。}
\cref{ch:clf-cbf} 把安全控制锻造成屏障函数 $h(x)$ 的前向不变性证书：只要 $\dot h\ge-\alpha(h)$，零超水平集 $\{h\ge0\}$ 就永不被离开。这套机制廉价、可在线求解 QP，但它把最难的一步——\emph{从哪里得到 $h$}——留给了工程师。这一步有三重困难。其一，\textbf{手工设计}：对低维、几何直观的约束（如“离障碍球心的距离”）可以凑出 $h$，但对高维、强耦合的动力学，一个既编码真实安全边界、又处处满足 CBF 条件的 $h$ 几乎无从构造。其二，\textbf{只看一步}：CBF 条件 $\dot h\ge-\alpha(h)$ 是\emph{瞬时}的、一步局部的判据，它保证“当前能不出界”，却不保证“从现在起始终有控制能不出界”——一个满足 CBF 条件的 $h$ 的安全集可能比真正的最大控制不变集大，落进去的状态最终仍会被逼出。其三，\textbf{扰动粗糙}：手工 $h$ 对有界扰动的处理通常靠加裕度（\cref{eq:issf-margin}），过紧则不安全、过松则过保守，无法精确刻画“最坏干扰下仍可控”的边界。

\paragraph{HJ 可达性的四个承诺。}
把安全验证反过来问——\textbf{有哪些初始状态，会在最坏情况干扰下不可避免地落入危险集？}——就把屏障函数的三重困难一举化解。这个“会出事的初始状态集合”就是后向可达管（BRT），它一旦算出，便给出四个承诺：(1) \textbf{自动}——$h$ 不再手设，而是从动力学 $f$ 与危险集 $g$ 经求解 HJI 方程算出；(2) \textbf{全局最优}——BRT 的补集恰是最大控制不变集（viability kernel），不存在更大的安全集，因此最不保守；(3) \textbf{最坏情况严格}——通过把有界扰动建模成对抗玩家并赋予其信息优势，得到的安全集是\emph{可证}的而非乐观的；(4) \textbf{附带最优控制}——$\arg\max_u$ 直接给出把系统留在安全集内的最小限制安全控制器（\cref{sec:dg-frs-cbf}）。代价是 HJI 需在状态空间网格上求解，遭遇维度灾难——这正是 CBF 廉价、HJ 精确的根本权衡。

\begin{table}[H]\centering\small
\caption{控制屏障函数（CBF）与 HJ 可达性的对照}
\label{tab:hj-cbf-vs-hj}
\begin{tabular}{lll}
\toprule
 & 控制屏障函数（\cref{ch:clf-cbf}） & HJ 可达性（本章） \\
\midrule
安全证书 $h/V$ 来源 & 手工设计 & 求解 HJI 方程算出 \\
判据时域 & 瞬时、一步局部 & 集合级、（无限）时域 \\
安全集 & 可能大于最大不变集（过乐观） & 精确的最大控制不变集 \\
最坏情况干扰 & 加裕度近似 & 对抗玩家、信息优势、精确 \\
计算时机/代价 & 在线 QP，便宜 & 离线 PDE，维度灾难 \\
适用维度 & 不限（受 $h$ 设计难度限） & 网格法 $n\lesssim6$，神经法更高但精度降 \\
\bottomrule
\end{tabular}
\end{table}

\begin{insight}
\textbf{“离线算清楚 + 在线快决策”是 HJ 与 CBF 的正确关系，而非二选一。}\;\cref{tab:hj-cbf-vs-hj} 看似把二者对立，但工程上最佳实践是把它们串起来：离线用 HJ 可达性把精确安全集与值函数 $V$ 算出来，再把 $V$ 当作一个\emph{现成的、最优的}屏障函数喂给在线 CBF-QP 执行。CBVF（\cref{sec:hj-cbvf}）正是把这条“离线 HJ $\to$ 在线 CBF”的接缝从理论上焊死的工具。换言之，HJ 不与 CBF 竞争，而是替 CBF 解决“$h$ 从哪来”这个它自己回答不了的问题。
\end{insight}

\paragraph{历史脉络。}
这条线索的根在 Isaacs 1965 年的《微分博弈》——他用追逃问题奠定了零和微分博弈与“Isaacs 方程”\cite{isaacs1965differential}。值函数在棱点不可微的困难，由 Crandall 与 Lions 1983 年的黏性解理论彻底解决\cite{crandall1983viscosity}。把演化界面表示成水平集、用 PDE 推进的数值框架，来自 Osher 与 Sethian 1988 年的奠基工作\cite{osher1988fronts}。Mitchell 2004 年发布首个开源水平集可达工具箱 toolboxLS（MATLAB）\footnote{I.\,M.\,Mitchell, \emph{A Toolbox of Level Set Methods}, UBC CS 技术报告 TR-2004-09（v1.0, 2004）、TR-2007-11（v1.1, 2007）。\url{https://www.cs.ubc.ca/~mitchell/ToolboxLS/}。}，Mitchell、Bayen 与 Tomlin 同年（2005 发表）给出连续动态博弈可达集的时变 Hamilton--Jacobi 表述，是工程领域的奠基论文（\emph{IEEE Trans.\ Automatic Control} 50(7):947--957）\cite{mitchell2005time}。此后 FaSTrack（Herbert 等，2017，CDC）把规划与安全解耦\cite{herbert2017fastrack}、DeepReach（Bansal 与 Tomlin，2017 提出、ICRA 2021 发表）用神经网络上探高维\cite{bansal2021deepreach}、CBVF（Choi 等，2021，CDC）把 CBF 与 HJ 统一\cite{choi2021cbvf}。最全面的入门综述是 Bansal、Chen、Herbert、Tomlin（2017，CDC）\cite{bansal2017hamilton}。

% ════════════════════════════════════════════════════════════════
\section{后向可达管的障碍算子：动态规划起源}
\label{sec:hj-brt-obstacle}

\cref{ch:plan-diffgame} 已给出 BRS、BRT、reach--avoid 三者的集合定义与对应 HJI 方程。这一节只做一件 \cref{ch:plan-diffgame} 提名而未展开的增量：把 BRT 方程里那个 $\min\{0,H\}$ 的截断项（\cref{eq:dg-brt-mitchell}）从动态规划原理\emph{推}出来，让它从“神秘的冻结项”变成“不可逆失败的记忆”，并补一个 Isaacs 不可分离的具体反例。

\subsection{量词结构与 Hamiltonian 中的 min/max}
\label{sec:hj-quantifier}

三类可达性问题的差别，全在“对控制问存在还是任意、对扰动问存在还是任意”的量词结构上，而量词结构唯一地决定 Hamiltonian 里对每个玩家取 $\min$ 还是 $\max$。这套对应（“存在取 $\min$、任意取 $\max$”）在 \cref{sec:dg-brs} 已讲透，这里仅以一表收束，供本章后续引用。

\begin{table}[H]\centering\small
\caption{可达性问题的量词结构与 Hamiltonian 取法（集合定义见 \cref{ch:plan-diffgame}）}
\label{tab:hj-quantifier}
\begin{tabular}{llll}
\toprule
问题 & 控制 $u$ & 扰动 $d$ & Hamiltonian \\
\midrule
BRS / BRT（不可避免落入危险） & $\forall u$（任意挡） & $\exists d$（能驱入） & $\max_u\min_d\lambda\cdot f$ \\
（控制方）可达目标 & $\exists u$（能到达） & $\forall d$（任意捣乱） & $\min_u\max_d\lambda\cdot f$ \\
reach--avoid（到达且全程避险） & $\exists u$（找到安全路径） & $\forall d$ & $\min_u\max_d$，外加双障碍 \\
\bottomrule
\end{tabular}
\end{table}

\noindent{}本章默认采用安全验证的角色：危险集为目标，控制力图\emph{远离}（故 $\forall u$ 取 $\max$）、扰动力图\emph{驱入}（故 $\exists d$ 取 $\min$），Hamiltonian 为 \cref{eq:dg-ham} 的 $H=\max_u\min_d\lambda\cdot f$。

\subsection{障碍算子从动态规划一步推出}
\label{sec:hj-dpp-obstacle}

BRT 回答的是“在时间窗 $[t,0]$ 内\emph{曾经}进入危险集”，而非“恰好在终点进入”。把这句话翻成值函数的递推，障碍算子就自然冒出来。在反向时间约定下（$t<0$ 表示距终端还有 $|t|$ 时长），令 $V(x,t)$ 表示从 $(x,t)$ 出发、在最优对抗下\emph{窗内最小的}带符号距离值——它 $\le0$ 当且仅当系统在窗内某刻进入了危险集。考虑一个微小时间步 $h>0$，把 $[t,0]$ 拆成“当前一刻”与“此后 $[t+h,0]$”两段，则动态规划原理给出

\begin{equation}
  V(x,t)=\max_{u(\cdot)}\min_{d(\cdot)}\ \min\Big\{\,\underbrace{g(x)}_{\text{此刻的危险距离}},\ \underbrace{V\big(\xi(t{+}h),\,t{+}h\big)}_{\text{此后窗内的最优值}}\,\Big\},
  \label{eq:hj-dpp-brt}
\end{equation}

其中 $\xi$ 是从 $x$ 出发、在 $u,d$ 作用下演化到 $t+h$ 的状态。\cref{eq:hj-dpp-brt} 外层的 $\min\{g(x),\cdot\}$ 是关键：它把“此刻已在危险集内”（$g(x)\le0$）这一事实\emph{永久记录}下来——一旦某刻 $g\le0$，外层 $\min$ 就把 $V$ 钉在 $\le0$，此后系统即便荡出危险集、$g$ 回升，这个 $V$ 也不会再回到正值。

\begin{derivation}[从 \cref{eq:hj-dpp-brt} 到障碍算子 PDE]
对 \cref{eq:hj-dpp-brt} 在 $h\to0$ 作一阶展开。在内层 $\min$ 取到第二支（即 $g(x)$ 不是最小，系统当前安全）时，标准 HJ 推导给出 $V(\xi(t{+}h),t{+}h)\approx V(x,t)+h\,[D_tV+\max_u\min_d\nabla_xV\cdot f]=V(x,t)+h\,[D_tV+H]$。代回并减去 $V(x,t)$、除以 $h$，得两支的竞争：

\[
0=\min\Big\{\,g(x)-V(x,t),\ \ D_tV+H(x,\nabla_xV)\,\Big\}.
\]

这正是 \cref{eq:dg-brt-vi} 的变分不等式形式。当目标用带符号距离编码、且 $V$ 已被初始化为 $V(x,0)=g(x)$ 时，可证它与把整个 Hamiltonian 截断的 Mitchell 形式

\begin{equation}
  D_tV+\min\big\{\,0,\ H(x,\nabla_xV)\,\big\}=0,\qquad V(x,0)=g(x)
  \label{eq:hj-brt-mitchell}
\end{equation}

等价（见 \cref{eq:dg-brt-mitchell}）：$\min\{0,H\}$ 中的“$0$”这一支阻止 $V$ 在已进入危险的区域随时间回升，效果与外层 $\min\{g-V,\cdot\}$ 的冻结一致。
\end{derivation}

\begin{insight}
\textbf{障碍算子 = 不可逆失败的记忆。}\;\cref{eq:hj-brt-mitchell} 里 $\min\{0,\cdot\}$ 之所以神秘，是因为孤立地看它只是个截断；放回 \cref{eq:hj-dpp-brt} 的动态规划语境就一目了然——它编码的是“撞过就算撞”这一物理事实。BRS（无此项，\cref{eq:dg-brs-pde}）只看终端一刻，会把“中途撞了、终点又恰好弹出危险区”的状态误判为安全；BRT（有此项）记住整个窗内的最坏时刻。安全验证几乎总要 BRT，正因为碰撞、坠落、越界这类事件都是不可逆的——发生过就计入，不因之后恢复而抹除。这也解释了为什么把 BRS 方程误用于安全分析会给出偏乐观、甚至致命的结论。
\end{insight}

\subsection{HJI 各项物理含义与 Isaacs 不可分离反例}
\label{sec:hj-isaacs}

\cref{eq:hj-brt-mitchell} 四个成分各有所指：$D_tV$ 是值随（反向）时间的变化——危险集如何向外“生长”成不安全集；$\nabla_xV\cdot f$ 是值沿最优对抗轨迹的方向导数——博弈双方如何沿动力学推拉值函数；$\min\{0,\cdot\}$ 是上述的不可逆记忆；终端条件 $V(x,0)=g(x)$ 把危险集播种为初值。其中 $\max_u\min_d$ 能写成对每个玩家独立优化、再合并的形式，依赖一个非平凡条件。

Isaacs 条件（鞍点存在，$\max_u\min_d=\min_d\max_u$，见 \cref{eq:dg-isaacs-cond}）保证上下值相等、博弈良定义。\cref{ch:plan-diffgame} 给出了它的一个充分条件——Hamiltonian 在 $(u,d)$ 上\emph{可分离}（控制项与扰动项可加性分开，\cref{eq:dg-separable}）。控制仿射、扰动仿射的系统 $f=f_0(x)+g_1(x)u+g_2(x)d$ 自动可分离，故工程中 Isaacs 条件几乎总成立。但可分离只是\emph{充分}条件，不是必要条件——下面这个不可分离的例子说明：动力学对 $(u,d)$ 乘性耦合，鞍点照样可以存在。

\begin{example}[不可分离但仍有鞍点：标量双线性 $\dot x=u\,d$]
\label{ex:hj-isaacs-fail}
考虑一维系统 $\dot x=u\cdot d$，控制 $u\in[-1,1]$、扰动 $d\in[-1,1]$，被积项 $p\cdot f=p\,u\,d$（$p=\nabla_xV$）。这是\emph{双线性}的——$u$ 与 $d$ 相乘，不可加性分离，故\cref{eq:dg-separable} 的可分离充分条件\emph{不}适用。直接算两个次序：

\[
\max_{u}\min_{d}\ p\,u\,d=\max_{u}\big(-|p\,u|\big)=\max_u(-|p|\,|u|)=0\quad(\text{取 }u=0),
\]
\[
\min_{d}\max_{u}\ p\,u\,d=\min_{d}\big(|p\,d|\big)=\min_d(|p|\,|d|)=0\quad(\text{取 }d=0).
\]

两次序都等于 $0$、相等：鞍点 $(u^\star,d^\star)=(0,0)$ 存在，Isaacs 条件成立。这\emph{不}是退化巧合——$pud$ 对固定 $d$ 关于 $u$ 仿射（故拟凹）、对固定 $u$ 关于 $d$ 仿射（故拟凸），在两个紧区间的乘积上由 von~Neumann--Sion 极小极大定理保证鞍点必存在。\footnote{标量双线性 $p\,u\,d$ 在任何轴对齐盒 $[u_1,u_2]\times[d_1,d_2]$ 上都满足 $\max_u\min_d=\min_d\max_u$：它对每个变量分别仿射、定义域为紧凸集，von~Neumann 极小极大定理（推广为 Sion 1958，\emph{Pacific J.\ Math.}\ 8:171--176）即给出鞍点。例如把控制集改为 $u\in[1,2]$（不含 $0$），逐点最优仍配平为 $\max_u\min_d=\min_d\max_u=-|p|$（取 $u^\star=1$、$d^\star=-\operatorname{sgn}p$），并\emph{不}破坏鞍点。故仅靠标量乘性耦合无法构造 Isaacs 失效。}

要真正让 Isaacs 失效（$\max_u\min_d<\min_d\max_u$ 严格），须破坏逐点 Hamiltonian 的“关于 $u$ 拟凹、关于 $d$ 拟凸”结构——这正是 Sion 定理的前提。注意\emph{任何}双线性耦合 $u^\top\mathbf M d$（无论 $\mathbf M$ 是否不定）都对每个玩家分别仿射，故在盒约束上仍满足前提、仍有鞍点；要破坏前提须让某一方\emph{非}仿射地进入。两条干净的失效路径：(i) 动作集\emph{离散}——如 \cref{ch:plan-diffgame} 的 $2\times2$ 矩阵博弈 $\big[\begin{smallmatrix}3&1\\2&4\end{smallmatrix}\big]$（无纯策略鞍点，上下值 $3>2$，须靠混合策略恢复）；(ii) Hamiltonian 对扰动\emph{非}拟凸——例如 $H(u,d)=-(u-d)^2$ 在 $u,d\in[-1,1]$ 上对 $d$ 是\emph{凹}的，逐点直接算得 $\max_u\min_d H=-1<0=\min_d\max_u H$，上下值严格不等。标量 $\dot x=ud$ 不属此列：其 Hamiltonian 双线性、对两方都仿射，鞍点保住。
\end{example}

\begin{remark}
\cref{ex:hj-isaacs-fail} 的教训是：可分离只是充分非必要——乘性耦合（乃至任意双线性耦合 $u^\top\mathbf M d$）\emph{不}自动导致鞍点消失。真正令 Isaacs 失效的，是逐点 Hamiltonian 失去“关于 $u$ 拟凹、关于 $d$ 拟凸”的 Sion 结构（离散动作集，或某一方非仿射地进入致拟凹/拟凸破坏，如对扰动凹的 $-(u-d)^2$）；此时 $\max_u\min_d$ 与 $\min_d\max_u$ 给出两个不同的值函数，分别对应“控制先动”与“扰动先动”的信息模式。安全验证取保守者（$\max\min$，给扰动信息优势，见 \cref{sec:dg-frs-cbf}）。值函数在棱点的不可微性以及“值函数 = 唯一黏性解”的严格基础见 \cref{thm:dg-evans}，本章不重复。
\end{remark}

% ════════════════════════════════════════════════════════════════
\section{水平集数值方法 I：迎风、Lax--Friedrichs 与数值黏性}
\label{sec:hj-numerics}

HJI 方程是连续状态空间上的一阶偏微分方程，须离散到网格上求解。\cref{ch:plan-diffgame} 已给出 Lax--Friedrichs（LF）数值 Hamiltonian 的公式（\cref{eq:dg-lf}）与 WENO/TVD-RK 的工程要点。本章补 \cref{ch:plan-diffgame} 未展开的双曲偏微分方程\emph{理论}视角：迎风差分“为什么必须”、LF 的人工黏性等价于数值扩散、以及 CFL 条件的严格推导。这些“为什么”是理解“单调格式为何收敛到黏性解”的钥匙。

\subsection{网格、反向时间步进与数值 Hamiltonian}
\label{sec:hj-grid}

把状态空间每维用 $N_i$ 个点离散，间距 $\Delta x_i$，总点数 $N=\prod_i N_i$，在每个网格点存一个 $V$ 值。从终端条件 $V(x,0)=g(x)$ 出发沿反向时间一步步推进：每步用相邻网格点的值近似空间梯度 $\nabla_xV$、算 Hamiltonian、再更新该点。难点全在“怎么近似 $\nabla_xV$”——因为值函数有不可微的棱（\cref{def:dg-viscosity}），天真的近似会在棱处崩掉。沿用 \cref{eq:dg-lf} 的 LF 数值 Hamiltonian：

\begin{equation}
  \hat H(x,p^+,p^-)=H\!\left(x,\frac{p^++p^-}{2}\right)-\sum_{i=1}^{n}\alpha_i\,\frac{p_i^+-p_i^-}{2},
  \label{eq:hj-lf}
\end{equation}

其中 $p_i^+=(V_{j+1}-V_j)/\Delta x_i$ 是第 $i$ 维前向差分、$p_i^-=(V_j-V_{j-1})/\Delta x_i$ 是后向差分，$\alpha_i=\max_{x,u,d}|\partial H/\partial p_i|$ 是第 $i$ 维数值黏性系数（即该维最大信息传播速度）。第一项取前后差分的平均作为梯度近似，第二项 $-\sum_i\alpha_i(p_i^+-p_i^-)/2$ 是人工黏性，正比于梯度跳变。工程上 $\alpha_i$ 逐网格点本地计算的\emph{局部} LF 是标准选择（全局 LF 取全网格最大值，简单但过耗散）。

\subsection{为什么必须迎风}
\label{sec:hj-upwind}

\paragraph{反面：中心差分的伪振荡。}
最自然的想法是用中心差分 $\partial_xV\approx(V_{j+1}-V_{j-1})/(2\Delta x)$。在光滑区它没问题，但在值函数的棱附近会产生剧烈伪振荡：它对称地混用棱两侧的信息，“看不见”间断的方向性，使数值解上下抖动、长出毛刺，可达集边界锯齿化甚至断裂。更糟的是，不带耗散的格式可能收敛到一个\emph{几乎处处满足方程、却不是黏性解}的假解（如 eikonal 方程的 $|x|-1$ 类被否决的解，见 \cref{sec:dg-viscosity}）。

\paragraph{双曲 PDE 的特征线决定信息方向。}
HJI 是一阶双曲型 PDE（不含二阶导），其解沿\textbf{特征线}传播信息。对一维线性输运方程 $V_t+aV_x=0$（$a>0$），信息从左向右传播——$(x,t)$ 处的值由上游点 $(x-a\Delta t,\,t-\Delta t)$ 决定。迎风差分的核心思想是\emph{只用信息流上游的差分}：$a>0$ 时用后向差分 $(V_j-V_{j-1})/\Delta x$（从左取信息），$a<0$ 时用前向差分。中心差分同时用了上下游，引入与物理传播方向相反的信息，这就是伪振荡的根源。HJ 方程尤其需要迎风，因为它的解常出现\textbf{扭折}（梯度不连续点），扭折处左右特征线汇聚成激波，不迎风就会在扭折附近野蛮振荡，彻底破坏零水平集精度。

\subsection{LF 隐式迎风等价于数值扩散}
\label{sec:hj-lf-diffusion}

LF 格式表面上是“中心差分 + 人工黏性”，看不出迎风。把 \cref{eq:hj-lf} 的黏性项按差分定义展开，迎风机制就显形了：

\begin{derivation}[LF 人工黏性的扩散形式]
将 $p_i^\pm$ 的定义代入 \cref{eq:hj-lf} 的黏性项：
\[
-\sum_i\frac{\alpha_i}{2}(p_i^+-p_i^-)
=-\sum_i\frac{\alpha_i}{2}\left(\frac{V_{j+1}-V_j}{\Delta x_i}-\frac{V_j-V_{j-1}}{\Delta x_i}\right)
=-\sum_i\frac{\alpha_i}{2\Delta x_i}\big(V_{j+1}-2V_j+V_{j-1}\big).
\]
括号里的 $V_{j+1}-2V_j+V_{j-1}$ 正是二阶中心差分（离散 Laplacian）的分子，故该项是一个\textbf{数值扩散}项，等效黏性系数为
\begin{equation}
  \varepsilon_i^{\mathrm{num}}=\frac{\alpha_i\,\Delta x_i}{2}.
  \label{eq:hj-numdiff}
\end{equation}
它在光滑处（$V_{j+1}-2V_j+V_{j-1}\approx0$）几乎消失，在扭折处（二阶差分大）自动提供稳定化。
\end{derivation}

\begin{insight}
\textbf{数值耗散 = 消失黏性的离散身影。}\;\cref{eq:hj-numdiff} 的人工黏性 $\alpha_i\Delta x_i/2$ 正是 \cref{ch:plan-diffgame} 黏性解理论里“消失黏性”$\varepsilon\Delta V$（$\varepsilon\to0^+$ 的极限挑出唯一黏性解，见 \cref{def:dg-viscosity}）的离散版本。理论用“加无穷小扩散再取极限”在数学上选解，数值用“加一点正比于网格的耗散”在计算上逼近同一个解：网格越细，等效 $\varepsilon_i^{\mathrm{num}}\propto\Delta x_i$ 越小，数值解越接近真黏性解。区别在于理论里 $\varepsilon\to0$ 是取极限，数值里耗散量由网格步长定、不能真取 $0$（否则振荡回来）。这条对应是“单调格式收敛到黏性解”的根据——收敛不是巧合，而是格式内建的耗散在兑现黏性极限。
\end{insight}

% ════════════════════════════════════════════════════════════════
\section{水平集数值方法 II：时间积分、CFL 与维度灾难}
\label{sec:hj-cfl}

空间离散后，HJI 变成一个大型常微分方程组 $\mathrm dV_j/\mathrm dt=-\min\{0,\hat H(x_j,p_j^+,p_j^-)\}$，须用时间积分推进。本节给时间格式、CFL 条件的严格推导与维度灾难的量化。

\subsection{TVD Runge--Kutta 时间积分}
\label{sec:hj-tvdrk}

为不引入伪振荡，时间积分用总变差不增（Total Variation Diminishing, TVD）的 Runge--Kutta 方法。常用三阶 TVD-RK3（Shu--Osher）：

\begin{equation}
\begin{aligned}
  V^{(1)}&=V^{n}+\Delta t\,L(V^{n}),\\
  V^{(2)}&=\tfrac{3}{4}V^{n}+\tfrac{1}{4}V^{(1)}+\tfrac{1}{4}\Delta t\,L(V^{(1)}),\\
  V^{n+1}&=\tfrac{1}{3}V^{n}+\tfrac{2}{3}V^{(2)}+\tfrac{2}{3}\Delta t\,L(V^{(2)}),
\end{aligned}
\label{eq:hj-tvdrk3}
\end{equation}

其中 $L(V)=-\min\{0,\hat H\}$ 是空间算子的离散。TVD-RK3 把三步的凸组合系数选得使每个中间级都是前向 Euler 步的凸组合，从而继承前向 Euler 的 TVD 性质——只要每个前向 Euler 子步满足 CFL，整体便不增总变差。

\subsection{CFL 条件的严格推导}
\label{sec:hj-cfl-derive}

CFL（Courant--Friedrichs--Lewy）条件不是经验法则，而有严格推导。从一维线性输运 $V_t+aV_x=0$（$a>0$）入手，用前向 Euler 时间 + 迎风（后向）空间：

\begin{equation}
  V_j^{n+1}=V_j^{n}-a\frac{\Delta t}{\Delta x}\big(V_j^{n}-V_{j-1}^{n}\big)=(1-\nu)V_j^{n}+\nu V_{j-1}^{n},\qquad \nu:=\frac{a\Delta t}{\Delta x},
  \label{eq:hj-cfl-1d}
\end{equation}

其中 $\nu$ 是 Courant 数。\cref{eq:hj-cfl-1d} 把 $V_j^{n+1}$ 写成 $V_j^n$ 与 $V_{j-1}^n$ 的组合，它是\emph{凸}组合当且仅当两系数非负，即 $0\le\nu\le1$。若 $\nu>1$，则 $1-\nu<0$，更新从“内插”退化为“外推”，误差以指数速度放大——物理含义是一个时间步内信息传播距离 $a\Delta t$ 超过了一个网格 $\Delta x$，信息“跳过”了网格点，格式无法捕捉。故一维 CFL 为 $\nu\le1$，即 $\Delta t\le\Delta x/a$。推广到多维 HJI，每维最大信息速度为 $\alpha_i=\max_{x,u,d}|f_i|$，由维度分裂的稳定性分析，各维 Courant 数之和不超过 1：

\begin{equation}
  \boxed{\ \sum_{i=1}^{n}\frac{\alpha_i\,\Delta t}{\Delta x_i}\le1\quad\Longleftrightarrow\quad \Delta t\le\frac{1}{\sum_{i=1}^{n}\alpha_i/\Delta x_i}.\ }
  \label{eq:hj-cfl-nd}
\end{equation}

实践中取 $\Delta t=\mathrm{CFL}_{\mathrm{safe}}\big/\sum_i(\alpha_i/\Delta x_i)$，$\mathrm{CFL}_{\mathrm{safe}}\approx0.8$ 留 $20\%$ 余量。对非线性系统 $\alpha_i$ 随 $V$ 演化而变，须每步重算并自适应调整 $\Delta t$。\cref{eq:hj-cfl-nd} 还揭示一个隐性代价：维度越高，$\sum_i\alpha_i/\Delta x_i$ 越大、允许的 $\Delta t$ 越小、总步数越多——这是高维 HJ 昂贵的原因之一。

\begin{pitfall}[编程陷阱：手动设过大的 $\Delta t$ 绕过 CFL]
为加速计算手动设 $\Delta t$（如 $0.1$）而不按 \cref{eq:hj-cfl-nd} 检查，是最常见的崩溃源。\textbf{现象}：值函数前几步看似正常，随后突然出现 \texttt{NaN} 或到处变成 $\pm10^{15}$，可达集形状完全错乱。\textbf{根因}：违反 CFL 后数值误差指数增长——以 $\Delta x\approx0.2$、$\alpha\approx5$ 的三维问题，\cref{eq:hj-cfl-nd} 要求 $\Delta t\le1/(3\times5/0.2)\approx0.013$，远小于 $0.1$。\textbf{自检}：算完检查 $|V|$ 是否合理，若 $|V|>100\,|g|_{\max}$ 几乎必是 CFL 违例。\textbf{正解}：用自动 CFL 计算，或手算并乘安全因子 $0.8$。另一个相关误区是“网格越细越好”：在 $n$ 维把每维加密 $2$ 倍，内存升 $2^n$、$\Delta t$ 减半使步数加倍、每步算量升 $2^n$，总算量升 $2^{n+1}$（4D 加密一次即 $32$ 倍），故应先用两个分辨率比较零水平集的 Hausdorff 距离判断是否已收敛，未收敛时优先提高空间格式阶数（如 WENO5）而非盲目加密。
\end{pitfall}

\subsection{维度灾难与工具谱系}
\label{sec:hj-curse}

网格法的致命弱点是维度灾难：存值函数需 $O(N^d)$ 内存（$d$ 为状态维、$N$ 为每维网格数）。\cref{tab:hj-curse} 把这一指数爆炸量化。

\begin{table}[H]\centering\small
\caption{水平集方法的内存随维度爆炸（float64，每点 8 字节）}
\label{tab:hj-curse}
\begin{tabular}{cccc}
\toprule
状态维 $d$ & 每维网格 $N_i$ & 总网格点 & 内存（float64） \\
\midrule
2 & 100 & $10^{4}$ & 80 KB \\
3 & 100 & $10^{6}$ & 8 MB \\
4 & 100 & $10^{8}$ & 800 MB \\
5 & 50 & $3\times10^{8}$ & 2.4 GB \\
6 & 40 & $4\times10^{9}$ & 32 GB \\
7 & 30 & $2\times10^{10}$ & 160 GB \\
\bottomrule
\end{tabular}
\end{table}

\noindent{}加上 CFL 对 $\Delta t$ 的压缩与每步 $O(N^d)$ 的算量，\textbf{网格法实际可算维度约 $d\le5\sim6$}。六维近悬停四旋翼已是网格法的极限，再高须靠分解（\cref{sec:hj-multiagent}）或神经逼近（\cref{sec:hj-deepreach}）。这一硬限正是 DeepReach 等方法出现的根本动机。

这些算法有成熟的开源实现，构成一条“网格法工具谱系”：toolboxLS（Mitchell，MATLAB，$d\le4$，经典标杆）与基于它、优化了内存的 helperOC（Berkeley，$d\le6$）；BEACLS（C++/CUDA，GPU 加速，$d\le6$）；以及更现代的、基于 JAX 自动微分与 JIT 的 \texttt{hj\_reachability}\footnote{\texttt{hj\_reachability}（Stanford ASL，JAX 实现，MIT 许可）：\url{https://github.com/StanfordASL/hj_reachability}。该仓库无正式配套论文，作者归属以仓库为准。}与基于 HeteroCL 的 \texttt{optimized\_dp}（SFU）。这些库的 API 细节属工程范畴，此处只取其谱系与适用维度；本章关心的是它们共同实现的数学——反向时间步进 + LF/WENO 数值 Hamiltonian + 障碍算子 + CFL 自适应步长。

\subsection{HJI 与 HJB 的关系}
\label{sec:hj-vs-hjb}

\cref{ch:dp-hjb} 建立了单智能体的 Hamilton--Jacobi--Bellman（HJB）方程与黏性解的一般理论。HJI 可视为 HJB 在三个方向的推广，三点对照即可（黏性解存在唯一性引 \cref{ch:dp-hjb} 与 \cref{thm:dg-evans}，本章不重证）。

\begin{table}[H]\centering\small
\caption{HJB（\cref{ch:dp-hjb}）与 HJI（本章/\cref{ch:plan-diffgame}）三点对照}
\label{tab:hj-hjb-hji}
\begin{tabular}{lll}
\toprule
 & HJB（单智能体最优控制） & HJI（零和博弈/可达性） \\
\midrule
对手 & 无 & 对抗扰动 $d$，Hamiltonian 含 $\min_d$ \\
代价 & 运行代价 + 终端代价 & 纯终端代价 $g$（去运行代价） \\
方程结构 & $D_tV+\min_u\{L+\nabla V\cdot f\}=0$ & $D_tV+\min\{0,\max_u\min_d\nabla V\cdot f\}=0$ \\
\bottomrule
\end{tabular}
\end{table}

\noindent{}三点合起来：HJI = HJB + 对手 $d$ + 去运行代价 + 障碍算子。去掉对手、加回运行代价，HJI 即退化为 HJB；这也说明可达性的数值机器（迎风、LF、CFL）可原样用于一般 HJB 黏性解的求解。

% ════════════════════════════════════════════════════════════════
\section{CBVF：控制屏障函数与 HJ 可达性的折扣统一}
\label{sec:hj-cbvf}

本节是全章第一理论高峰，也是 \cref{ch:clf-cbf} 明确“给桥、I 章给主证”的去重接口（\cref{eq:cbvf} 在 F 章给出 CBVF 的桥接陈述并指向本节主证）。它把看似对立的两条路——\cref{sec:dg-frs-cbf} 末尾点出的“HJ 精确但贵 vs CBF 廉价但保守”——用一个折扣参数 $\gamma$ 连成一条连续谱。控制屏障-值函数（Control Barrier-Value Function, CBVF）由 Choi、Lee、Sreenath、Tomlin、Herbert（2021，CDC）提出\cite{choi2021cbvf}。

\subsection{折扣 HJI 变分不等式}
\label{sec:hj-cbvf-vi}

引入折扣 $\gamma\ge0$，CBVF $B_\gamma(x,t)$ 定义为如下折扣 HJI 变分不等式的黏性解（沿用 Choi 2021 的标准 VI 形式，与 \cref{eq:cbvf} 一致；终端 $B_\gamma(x,0)=g(x)$，安全集取零超水平集 $\{B_\gamma\ge0\}$）：

\begin{equation}
  \boxed{\ 0=\min\Big\{\,g(x)-B_\gamma,\ \ D_tB_\gamma+\max_{u\in\mathcal U}\min_{d\in\mathcal D}\nabla_xB_\gamma\cdot f(x,u,d)+\gamma\,B_\gamma\,\Big\}.\ }
  \label{eq:hj-cbvf}
\end{equation}

与 BRT 的 \cref{eq:dg-brt-vi} 相比，唯一的新增是 Hamiltonian 一支里多出的 $+\gamma B_\gamma$ 项——它正是 CBF 衰减条件 $\dot h\ge-\alpha(h)$ 在 $\alpha(h)=\gamma h$（线性类 $\mathcal K$ 函数，见 \cref{def:ext-kinf}）时的内嵌。

\begin{note}
\textbf{两种障碍写法的等价（\cref{ch:plan-diffgame} 口径对齐）.}\;源文献亦常把 CBVF 写成 Mitchell 截断形式 $D_tB_\gamma+\min\{0,\,\max_u\min_d\nabla B_\gamma\cdot f+\gamma B_\gamma\}=0$。在标准 BRT 情形（终端 $B_\gamma(\cdot,0)=g$、目标用带符号距离编码）下，这一截断写法与 \cref{eq:hj-cbvf} 的变分不等式形式等价；本章以 Choi 2021 原文的 VI 形式 \cref{eq:hj-cbvf} 为准（与 \cref{eq:cbvf} 一致），障碍项写作 $g-B_\gamma$。符号约定也与原文一致：安全集 $\{B_\gamma\ge0\}$、控制取 $\max$、扰动取 $\min$。
\end{note}

\subsection{两个极端与连续谱}
\label{sec:hj-cbvf-spectrum}

\cref{eq:hj-cbvf} 的威力在于：$\gamma$ 从 $0$ 到 $\infty$ 扫过时，CBVF 从“经典 HJ 可达”连续变形到“纯 CBF”。

\begin{proposition}[$\gamma=0$ 恢复经典 HJ 可达]
\label{prop:hj-cbvf-zero}
当 $\gamma=0$ 时，\cref{eq:hj-cbvf} 退化为标准 BRT 变分不等式 \cref{eq:dg-brt-vi}，$\{B_0\ge0\}$ 是\emph{最大控制不变集}（viability kernel），不存在更大的安全集。
\end{proposition}

\begin{proof}
代 $\gamma=0$，\cref{eq:hj-cbvf} 中 $+\gamma B_\gamma$ 项消失，得 $0=\min\{g-B_0,\ D_tB_0+\max_u\min_d\nabla_xB_0\cdot f\}$，即 \cref{eq:dg-brt-vi}。$\gamma=0$ 不施加任何额外衰减约束，值函数的演化完全由动力学与障碍算子决定，故零超水平集是不可被改进的最大不变集（精确性见 \cref{ch:plan-diffgame}）。
\end{proof}

\begin{proposition}[$\gamma\to\infty$ 恢复 CBF]
\label{prop:hj-cbvf-inf}
当 $\gamma\to\infty$ 时，CBVF 的稳态解趋于满足 CBF 条件 $\dot B\ge-\gamma B$ 的屏障函数，安全集 $\{B_\gamma\ge0\}$ 退缩到由终端 $g$ 的原始零超水平集 $\{g\ge0\}$ 附近。
\end{proposition}

\begin{proof}[证明骨架]
在零超水平集内部（$B_\gamma>0$，障碍算子不起作用），稳态（$D_tB_\gamma=0$）方程为 $\max_u\min_d\nabla_xB_\gamma\cdot f+\gamma B_\gamma=0$，即 $\max_u\min_d\nabla_xB_\gamma\cdot f=-\gamma B_\gamma$。作重标度 $B_\gamma=W/\gamma$，方程化为 $\max_u\min_d\nabla_xW\cdot f/\gamma+W=0$；当 $\gamma\to\infty$，第一项 $\to0$，故 $W\to0$ 处处，即 $B_\gamma\to0$，安全集退缩到 $\{g\ge0\}$。同时在 $\{B_\gamma>0\}$ 内，稳态条件 $\max_u\min_d\nabla_xB_\gamma\cdot f\ge-\gamma B_\gamma$ 恰是 $\alpha(h)=\gamma h$ 对应的 CBF 不等式——$B_\gamma$ 本身就是一个合法的 CBF（前向不变性证明见 \cref{thm:cbf-main}）。
\end{proof}

\begin{theorem}[CBVF 三性质，Choi 2021 Thm.\ 1]
\label{thm:hj-cbvf}
对所有 $\gamma\ge0$，\cref{eq:hj-cbvf} 的黏性解 $B_\gamma$ 满足：
\begin{enumerate}
  \item（\textbf{单调性}）$B_\gamma(x,t)$ 关于 $\gamma$ 单调递增，故安全集 $\{B_\gamma\ge0\}$ 关于 $\gamma$ 单调\emph{递减}（$\gamma$ 越大越保守）；
  \item（\textbf{前向不变性}）$\{B_\gamma\ge0\}$ 是控制不变集——存在控制使系统不离开它；
  \item（\textbf{合法 CBF}）$B_\gamma$ 在 $\{B_\gamma\ge0\}$ 上自动满足 CBF 条件，可直接作为 CBF-QP 的屏障函数。
\end{enumerate}
\end{theorem}

\begin{proof}[证明思路]
性质 1 由比较原理：$+\gamma B$ 项使 Hamiltonian 一支随 $\gamma$ 增大，作为 VI 的解 $B_\gamma$ 随之逐点不减；超水平集随之收缩。性质 2 源于障碍项 $g-B_\gamma$ 把 $B_\gamma$ 封顶在 $g$ 之下并冻结，使 $\{B_\gamma\ge0\}$ 只增不减、最终成为不变集；这是 BRT 不变性（\cref{ch:plan-diffgame}）在折扣下的延拓。性质 3 由 \cref{prop:hj-cbvf-inf} 证明中导出的稳态不等式 $\max_u\min_d\nabla_xB_\gamma\cdot f\ge-\gamma B_\gamma$ 给出——这正是 \cref{def:zcbf} 形式的 CBF 条件。三性质合起来正是 CBVF 名称的由来：它既是博弈值函数（Value Function），又是屏障函数（Barrier Function）。
\end{proof}

\begin{table}[H]\centering\small
\caption{CBVF 连续谱：折扣 $\gamma$ 调节安全集与计算代价}
\label{tab:hj-cbvf-spectrum}
\begin{tabular}{llll}
\toprule
$\gamma$ & 行为 & 安全集 & 计算 \\
\midrule
$0$ & 经典 HJ 可达 & 最大（精确 BRT/viability kernel） & 需收敛到稳态，慢 \\
小 & 接近 HJ & 接近最大 & 收敛较快 \\
大 & 接近 CBF & 较小、保守 & 收敛很快 \\
$\infty$ & 纯 CBF & 取决于 $g$ 设计 & 无需离线计算 \\
\bottomrule
\end{tabular}
\end{table}

\begin{remark}
工程意义直白：过去工程师在“花大量离线算力得精确 BRT（$\gamma=0$）”与“用手设、可能保守的 CBF（$\gamma\to\infty$）”之间二选一；CBVF 给出中间档——用适当的 $\gamma$ 快速得到“够好”的安全控制器，既不像纯 HJ 那么贵，也不像纯 CBF 那么粗糙。实践上离线用 HJ 方法算 $B_\gamma$、在线用鲁棒 CBVF-QP 执行，从根上解决了“屏障函数从哪来”的痛点；唯一限制仍是 HJ 只能算低维（$n\lesssim6$）。这正是 \cref{eq:cbvf} 在 \cref{ch:clf-cbf} 所引的“HJ 给出最优 CBF”那句话的完整理论。
\end{remark}

\subsection{一维连续谱解析算例}
\label{sec:hj-cbvf-1d}

用一维系统看清 $\gamma$ 如何不动安全集边界、却改变值函数的衰减率。设 $\dot x=u+d$，$u\in[-1,1]$、$d\in[-0.5,0.5]$，危险目标 $\mathcal T=\{x\ge3\}$，故 $g(x)=3-x$（$g\le0$ 在危险集内）。最坏净速度 $=-1+0.5=-0.5$（控制总能克服扰动以 $0.5$ 远离危险）。

\begin{itemize}
  \item \textbf{$\gamma=0$（经典 HJ）}：从安全侧任一点，控制能以净速度 $0.5$ 远离危险集，故无限时域 BRT $=\{x\ge3\}=\mathcal T$ 本身，安全集 $=\{x<3\}$。
  \item \textbf{$\gamma=0.5$（中等折扣）}：在安全侧 $x<3$，$B>0$ 且 $\nabla B=B'<0$（越靠近危险越小），最优控制 $u=-1$、最坏扰动 $d=0.5$，稳态方程 $B'\cdot(-0.5)+0.5\,B=0$ 即 $B'=B$，解形如 $B(x)=C(\mathrm e^{x}-\mathrm e^{3})$ 满足边界 $B(3)=0$。关键结论：安全集边界\emph{仍在 $x=3$}（因边界上 $\gamma B=0$），但 $B_\gamma$ 远离边界时增长更慢——$\gamma$ 越大，$B_\gamma$ 在内部越“扁”，CBF-QP 用它时安全滤波器介入得越早、越保守。
  \item \textbf{$\gamma\to\infty$（CBF 极限）}：$B_\gamma\to g=3-x$，安全集 $=\{3-x\ge0\}=\{x\le3\}$，与 $g$ 的原始零超水平集重合，CBF 条件成 $-(u+d)\ge-\gamma(3-x)$。
\end{itemize}

\noindent{}三档安全集边界都在 $x=3$，区别只在值函数的形状（衰减率）——这把 \cref{thm:hj-cbvf} 的单调性在最简单的情形上具体化了。

\begin{remark}
不要以为“$\gamma$ 越小越好”。$\gamma=0$ 确实给最大安全集，但有两个代价：一是须完全收敛到稳态才准确、计算时间长；二是在高维（DeepReach，\cref{sec:hj-deepreach}）里 $\gamma>0$ 的衰减效应充当正则化、使训练更稳。Choi 2021 的实验表明 $\gamma\in[0.1,1.0]$ 通常是安全集大小与计算效率的良好折中。
\end{remark}

% ════════════════════════════════════════════════════════════════
\section{DeepReach：高维 HJ 的神经网络逼近}
\label{sec:hj-deepreach}

网格法在 $d>6$ 失效（\cref{sec:hj-curse}）。本节是全章第二理论增量，给 \cref{ch:plan-diffgame} 仅提名的 DeepReach（Bansal 与 Tomlin，ICRA 2021）一个深推导\cite{bansal2021deepreach}：用物理信息神经网络（PINN）参数化值函数 $V_\theta(x,t)$，把维度灾难从“内存”转移到“训练”。它接 RL 部的神经网络与自动微分基础（\cref{ch:rl-basic}），此处只用其结论。

\subsection{从 PINN 到 SIREN 架构}
\label{sec:hj-pinn}

PINN 的思路是把 PDE 残差当作损失：不依赖网格，而用神经网络 $V_\theta(x,t)$ 直接拟合，损失项是 HJI 方程在大量随机采样点上的残差加上终端条件残差。DeepReach 的关键架构选择是 SIREN（正弦表示网络，Sitzmann 等 2020\cite{sitzmann2020siren}）——每层用正弦激活：

\begin{equation}
  z^{(k+1)}=\sin\!\big(\omega_0\,(W^{(k)}z^{(k)}+b^{(k)})\big),
  \label{eq:hj-siren}
\end{equation}

频率超参 $\omega_0$ 控制网络能表示的频率范围（DeepReach 默认 $\omega_0=30$）。SIREN 配特定初始化（第一层权重 $\sim U(-1/n,1/n)$、后续层 $\sim U(-\sqrt{6/n}/\omega_0,\,\sqrt{6/n}/\omega_0)$，$n$ 为输入维），使各层激活分布在训练中稳定。

\begin{insight}
\textbf{为什么 SIREN 用 $\sin$ 而非 ReLU——HJI 需要 $\nabla V$ 这一事实逼出来的。}\;HJI 残差含 $\nabla_xV$，PINN 要最小化残差还须对残差再求导（即用到 $V$ 的二阶导）。ReLU 网络 $\sigma(x)=\max(0,x)$ 的二阶导几乎处处为零、一阶导分段常数，故 $\nabla V_\theta$ 在每个线性区内\emph{恒定}、梯度信号极稀疏——要逼近光滑弯曲的零水平集需海量折线段，效率极低。$\sin$ 激活则有三大优势：(1) \textbf{各阶导非零}——$\sin$ 的导数是 $\cos$、二阶是 $-\sin$、永不为零，PINN 的各阶导训练信号都丰富；(2) \textbf{频率可控}——$\omega_0$ 决定可表示频率，HJI 值函数在零水平集附近从正剧变到负、需要高频表示能力；(3) \textbf{初始化有数学保证}——Sitzmann 等证明上述初始化下网络初始输出服从反正弦分布、激活在训练中保持稳定，类似 He/Xavier 对 ReLU/Sigmoid 的作用。$\omega_0$ 过大会使损失景观崎岖、训练不稳，通常从 $30$ 起调：过平滑则增大、不稳则减小。
\end{insight}

\subsection{训练目标与课程学习}
\label{sec:hj-curriculum}

DeepReach 的损失是 PDE 残差与边界（终端）条件残差之和：

\begin{equation}
  \mathcal L(\theta)=\underbrace{\big\|\,D_tV_\theta+\min\{0,H(x,\nabla_xV_\theta)\}\,\big\|_{\Omega\times[t,0]}^{2}}_{\text{HJI 残差（在随机配点上）}}+\ \lambda\,\underbrace{\big\|\,V_\theta(x,0)-g(x)\,\big\|_{\Omega}^{2}}_{\text{终端条件残差}},
  \label{eq:hj-deepreach-loss}
\end{equation}

其中 $\lambda$（BC 权重，常取 $100$）平衡两项。直接在整个时间范围上训练很难收敛，DeepReach 用\textbf{课程学习}：先只在终端附近的小时间窗 $[-\epsilon,0]$ 上训练（此时 $V_\theta\approx g$，易学），再逐步把窗口向 $t$ 负方向扩展。这恰好\emph{模仿了水平集方法的反向时间步进}——先学好终端、再沿时间倒推，把一个难的全局拟合分解成一串由易到难的子问题。

\begin{table}[H]\centering\small
\caption{DeepReach 关键超参与调优方向（经验默认，非理论最优）}
\label{tab:hj-deepreach-hyper}
\begin{tabular}{lll}
\toprule
超参 & 推荐值 & 调优方向 \\
\midrule
隐藏层数 / 宽度 & 3 / 512 & $d>8$ 增至 4--5 层 / 复杂系统增至 1024 \\
$\omega_0$ & 30 & 见 \cref{sec:hj-pinn} \\
学习率 & $5\times10^{-5}$ & 训练不稳时减小 \\
BC 权重 $\lambda$ & 100 & 终端条件不满足时增至 500 \\
课程阶段数 & 10--20 & 时间范围大时增加 \\
\bottomrule
\end{tabular}
\end{table}

训练诊断看四个量：PDE 损失应单调下降（大幅震荡 $\Rightarrow$ 学习率或 $\omega_0$ 过大）；BC 损失应早期快速降到 $<10^{-4}$（居高 $\Rightarrow$ 增大 $\lambda$）；定期画 $V_\theta=0$ 等值线看零水平集是否渐光滑合理；低维时用网格法算 ground truth，测零水平集的 Hausdorff 距离与 IoU\footnote{实现上须对自动微分保留计算图（求 HJI 残差再反传需二阶导），并令输入 $x,t$ 可求导，否则 PDE 损失不下降——这是 PINN 训练的常见陷阱，但属框架细节，不展开。}。

\begin{remark}
\textbf{DeepReach 没有“解决”维度灾难，只是搬了家。}\;无网格不等于无诅咒——它把维度灾难从内存转移到了\emph{训练难度与采样稀疏}。高维空间随机采样极稀疏：$10$ 维里 $10^5$ 个均匀点的平均间距约 $10$ 个单位长度，值函数的零水平集（一张 $9$ 维超曲面）在采样点中被“看到”的概率极低，故 DeepReach 在 $10$ 维以上精度往往很差。它也没有网格法的收敛保证，且在低维上通常\emph{不如}网格法精确。诚实的定位是：DeepReach 是网格法力不能及的高维区的可用近似，不是免费午餐。前沿应对包括在零水平集附近自适应集中采样、把 HJI 结构编码进网络架构、以及低维网格 + 高维神经的混合分解。
\end{remark}

% ════════════════════════════════════════════════════════════════
\section{解耦式安全框架的理论}
\label{sec:hj-frameworks}

可达性的产物不止一个集合，还附带一个最优安全控制器，这使它能嵌入实际系统的安全架构。本节给三类解耦框架的完整数学，与规划部 \cref{sec:pr-reach}（已给 FaSTrack/RTD/Funnel 的规划用途与跟踪误差界直觉）互补去重：规划侧讲“换规划器不影响安全”的架构动机，控制侧（本节）讲“跟踪误差界怎么从 HJ 算出来”。

\subsection{FaSTrack：规划与安全解耦的完整数学}
\label{sec:hj-fastrack}

FaSTrack（Fast and Safe Tracking，Herbert 等 2017，CDC\cite{herbert2017fastrack}）把高维系统的安全保证拆成“低维快速规划 + 离线一次性安全计算”。\cref{sec:pr-reach} 已给其架构与用途，这里推完 TEB 的数学。

设\textbf{规划模型} $\dot s=g(s,u_p)$ 维度低（如 2D/3D 单积分器，规划器在其上用 A*/RRT 快速规划），\textbf{跟踪模型} $\dot x=f(x,u_t,d)$ 维度高（如 6D/10D 四旋翼，含扰动 $d$），跟踪误差 $e=\phi(x)-s$（$\phi$ 是高维状态到低维规划空间的投影）。误差动力学为

\begin{equation}
  \dot e=\frac{\partial\phi}{\partial x}\,f(x,u_t,d)-g(s,u_p).
  \label{eq:hj-teb-edot}
\end{equation}

FaSTrack 的关键一招：把规划器输出 $u_p$ 视为\textbf{对抗扰动}——规划器可任意改方向，跟踪器须在最坏情况下保证误差有界。于是对 \cref{eq:hj-teb-edot} 做 HJ 可达性分析，控制是跟踪控制 $u_t$（最小化误差）、扰动是 $(d,u_p)$ 联合体（最大化误差）、目标集是误差球 $\{e:\|e\|\le r\}$。算出 BRT 后，使最大控制不变集 $\mathcal I(r)$ 仍含于误差球的最大 $r$ 就是\textbf{跟踪误差界}（Tracking Error Bound, TEB）：

\begin{equation}
  \mathrm{TEB}=\max\big\{\,r:\ \mathrm{BRT}_\infty(r)\subseteq\{e:\|e\|\le r\}\,\big\}.
  \label{eq:hj-teb}
\end{equation}

实操上固定小 $r$ 算 BRT、若含于误差球则有效、逐步增大找最大值。有了 TEB，规划层把障碍按 Minkowski 和膨胀 $\mathcal O_{\mathrm{inflated}}=\mathcal O\oplus B(\mathrm{TEB})$（$B(\mathrm{TEB})$ 是半径 TEB 的球），只要低维规划路径不进膨胀障碍，跟踪器就保证高维真实系统不碰原始障碍。在线时跟踪控制器用误差空间的 HJ 值函数取 $u_t^\ast=\arg\max_{u_t}\min_{(d,u_p)}\nabla V_e\cdot\dot e$，保证误差永不超过 TEB。其妙处在于：规划器可任意替换而不影响安全保证（安全只依赖误差动力学的离线 BRT），且规划在低维空间、安全在误差空间，都回避了高维直接 HJ。

\begin{remark}
TEB 不是“误差为零”，而是误差\emph{有界} $\|e\|\le\mathrm{TEB}$。TEB 的大小取决于规划模型与跟踪模型能力的差距——差距越大，TEB 越大、障碍膨胀越多、可行路径空间越窄。设计指南：规划模型应尽量贴近真实系统能力（真实最大速度 $5$ m/s，规划模型也限 $5$ m/s 以内），否则 TEB 会很大。
\end{remark}

\subsection{Fisac 最小干预安全框架}
\label{sec:hj-fisac}

在人-机器人共享控制（外骨骼辅助、辅助驾驶）中，人类提供输入 $u_h$ 但可能犯错，如何在\emph{最小干预}人类意图的前提下保证安全？Fisac 等（2019，\emph{IEEE Trans.\ Automatic Control}\cite{fisac2019general}）把它建模为 $\dot x=f(x,u_h,u_s)$（$u_h$ 视为不确定/对抗，$u_s$ 是安全系统修正），用 HJ 值函数 $V$（$\{V\le0\}$ 为不安全的 BRT、$\{V>0\}$ 安全）做三档切换：

\begin{equation}
  u_s=\begin{cases}
    0, & V(x)>\epsilon\quad(\text{正常：远离边界，完全信任人类}),\\[2pt]
    \big(1-V(x)/\epsilon\big)\,u_{\mathrm{safe}}, & 0<V(x)\le\epsilon\quad(\text{警戒：渐进介入}),\\[2pt]
    u_{\mathrm{safe}}, & V(x)\le0\quad(\text{危急：完全接管}),
  \end{cases}
  \label{eq:hj-fisac}
\end{equation}

其中 $u_{\mathrm{safe}}=\arg\max_u\min_d\nabla V\cdot f$ 是 HJ 算出的最优安全控制。与 CBF-QP 相比：CBF-QP 在每个时刻都修正控制（可能频繁微调），Fisac 框架只在逼近 BRT 边界时介入（干预更少但更狠），这正是“最小限制”性质（\cref{sec:dg-frs-cbf}）在共享控制里的体现——避免传统安全系统过于保守地频繁打断人类。

\subsection{多智能体分解与投影}
\label{sec:hj-multiagent}

$N$ 个智能体（各 3D 状态）合起来是 $3N$ 维，直接 HJ 不可行。Chen、Herbert 等（2016/2018 系统分解\cite{chen2018decomposition}）的策略是\emph{两两分解}：对每对智能体 $(i,j)$ 在相对坐标下做 3D 的 HJ 分析，得两两安全控制策略 $V_{ij}$；当多个两两约束同时激活，用优先级排序裁决；智能体 $i$ 的安全控制取所有两两策略的最保守交集 $\arg\max_u\min_j V_{ij}(x_{\mathrm{rel}})$，共 $\binom N2$ 个 3D 计算替代一个 $3N$ 维计算。代价是保守——两两分解忽略三体及以上的联合交互（智能体 1 为避让 2 可能被迫靠近 3），真正的多体安全需联合可达，但目前仅限 2--3 个简单智能体。一种相关的降维是\textbf{投影}：对一个高维 BRT 沿某些维（如速度维）取 $\min$，得到低维的保守过近似 $\mathrm{BRT}_{\mathrm{proj}}$，可在线快速查询。这条多智能体线索与规划部的实时博弈、安全 MARL、机会约束求解器互补，本章不涉其求解细节，详见 \cref{ch:plan-rtgame,ch:plan-safemarl,ch:plan-chance}。工业上 ACAS X 空中防撞系统即把离线算好的安全策略压成查找表、在线按当前相对状态查表决策。

% ════════════════════════════════════════════════════════════════
\section*{本章常见误解汇总}

\begin{itemize}
  \item \textbf{“BRT 收敛了就该停止计算。”}\;BRT 随反向时间单调增长、最终收敛到无限时域可达集；但收敛是渐近的，过早停止会低估不安全集、高估安全集。判收敛应比较相邻时间步零水平集的变化量低于阈值，而非凭计算时长。
  \item \textbf{“把 BRS 方程（无障碍算子）用于安全验证。”}\;BRS 只看终端一刻，会漏掉“中途撞了、终点又弹出”的状态（\cref{ex:dg-set-vs-tube}），给出偏乐观甚至致命的结论。安全验证一律用带障碍算子的 BRT（\cref{eq:hj-brt-mitchell}）。
  \item \textbf{“$\gamma$ 越小越好。”}\;$\gamma=0$ 给最大安全集但须完全收敛、计算慢，且在高维神经逼近里失去正则化。$\gamma\in[0.1,1.0]$ 常是安全集大小与效率的折中（\cref{sec:hj-cbvf-1d}）。
  \item \textbf{“FaSTrack 的 TEB 保证零误差。”}\;TEB 是误差\emph{有界}，非零；其大小由规划模型与跟踪模型的能力差距决定（\cref{sec:hj-fastrack}）。
  \item \textbf{“DeepReach 解决了维度灾难。”}\;它把诅咒从内存转移到训练难度与采样稀疏，$10$ 维以上精度差、无收敛保证（\cref{sec:hj-deepreach}）。
  \item \textbf{“Isaacs 条件总是成立。”}\;控制/扰动仿射时自动可分离、Isaacs 成立；但 Isaacs 可在逐点 Hamiltonian 失去“关于 $u$ 拟凹、关于 $d$ 拟凸”结构时失效（离散动作集如矩阵博弈，或某一方非仿射进入），此时 $\max\min<\min\max$。\textbf{注意}：标量乘性耦合 $\dot x=u\,d$（乃至任意双线性 $u^\top\mathbf M d$）不足以破坏鞍点——它对两方分别仿射，由极小极大定理在盒约束上仍有鞍点（\cref{ex:hj-isaacs-fail}）。
\end{itemize}

% ════════════════════════════════════════════════════════════════
\section*{本章小结}

本章把 HJ 可达性从“集合级、最坏情况、无限时域安全验证”的角度讲透，定位为 \cref{ch:plan-diffgame} 奠基理论之上的四块控制部深化。其一，\textbf{障碍算子的动态规划起源}（\cref{sec:hj-brt-obstacle}）：$\min\{0,H\}$ 不是神秘的截断，而是“不可逆失败的记忆”，从 \cref{eq:hj-dpp-brt} 的动态规划原理一步推出；并以 $\dot x=u\,d$ 反例说明 Isaacs 可分离只是充分非必要。其二，\textbf{水平集数值方法的双曲 PDE 理论}（\cref{sec:hj-numerics,sec:hj-cfl}）：迎风差分由特征线方向\emph{必须}采用，LF 的人工黏性 $\alpha_i\Delta x_i/2$ 是消失黏性的离散身影（\cref{eq:hj-numdiff}），CFL 条件 \cref{eq:hj-cfl-nd} 由凸组合严格推出，维度灾难 $O(N^d)$ 把网格法压在 $d\lesssim6$。其三，\textbf{CBVF}（\cref{sec:hj-cbvf}）：折扣 HJI 变分不等式 \cref{eq:hj-cbvf} 用 $\gamma$ 把经典 HJ（$\gamma=0$，最大不变集）与纯 CBF（$\gamma\to\infty$）连成连续谱，三性质（\cref{thm:hj-cbvf}）使 $B_\gamma$ 既是值函数又是合法屏障函数——这是 \cref{ch:clf-cbf} 所引“HJ 给出最优 CBF”的主证。其四，\textbf{DeepReach 与解耦框架}（\cref{sec:hj-deepreach,sec:hj-frameworks}）：SIREN 用正弦激活以提供 PINN 所需的各阶导信号，课程学习模仿反向时间步进；FaSTrack 把规划器输出当对抗扰动、用误差动力学 BRT 算出 TEB，Fisac 三档最小干预切换，多智能体两两分解避开 $3N$ 维。

本质洞察是：HJ 可达性把“安全证书从哪来”这个 CBF 自己回答不了的问题，转化为一个可\emph{算}的最优控制问题——代价是维度灾难，回报是最不保守、可证、附带最优控制器的安全集。

\begin{table}[H]\centering\small
\caption{本章关键结论速查}
\label{tab:hj-summary}
\begin{tabular}{lll}
\toprule
主题 & 结论 & 出处 \\
\midrule
BRT 障碍算子 & $D_tV+\min\{0,H\}=0$，不可逆失败记忆 & \cref{eq:hj-brt-mitchell} \\
Isaacs 可分离性 & $\dot x=u\,d$ 不可分离却仍有鞍点（充分非必要） & \cref{ex:hj-isaacs-fail} \\
LF 数值黏性 & $\varepsilon_i^{\mathrm{num}}=\alpha_i\Delta x_i/2$ & \cref{eq:hj-numdiff} \\
多维 CFL & $\sum_i\alpha_i\Delta t/\Delta x_i\le1$ & \cref{eq:hj-cfl-nd} \\
维度灾难 & $O(N^d)$，网格法 $d\lesssim6$ & \cref{tab:hj-curse} \\
CBVF 变分不等式 & $\min\{g-B_\gamma,\ D_tB_\gamma+H+\gamma B_\gamma\}=0$ & \cref{eq:hj-cbvf} \\
CBVF 三性质 & 单调 / 前向不变 / 合法 CBF & \cref{thm:hj-cbvf} \\
FaSTrack TEB & 误差动力学 BRT 的水平集，障碍 Minkowski 膨胀 & \cref{eq:hj-teb} \\
\bottomrule
\end{tabular}
\end{table}

% ════════════════════════════════════════════════════════════════
\section*{通向高级：本章交给后续什么}

本章向四个方向延伸。其一，向 \cref{ch:clf-cbf}：CBVF（\cref{eq:hj-cbvf}）的三性质主证在此，F 章 \cref{eq:cbvf} 引用之，把“离线 HJ 算最优屏障函数、在线 CBF-QP 执行”这条线焊死——本章给主证、F 章给桥与 QP 综合。其二，向 \cref{ch:dp-hjb}：HJI = HJB + 对手 + 去运行代价 + 障碍算子（\cref{tab:hj-hjb-hji}），可达性的迎风/LF/CFL 数值机器可原样用于一般 HJB 黏性解求解，二者值函数视角同源。其三，向 \cref{ch:mpc-advanced}：MPC 安全滤波与 HJ 离线值函数互补——离线 HJ 把精确安全集与 $V$ 算清，在线把 $V$ 当屏障函数喂给 MPC/CBF-QP 快速决策，集合级的最坏情况保证补上单轨迹预测的盲点。其四，与 \cref{ch:ddp-ilqr} 形成路线对照：DDP/iLQR 是\emph{局部}轨迹优化（单条轨迹、有限时域、二次近似），HJ 是\emph{全局}可达（集合级、无限时域、PDE）——路线不同、互不替代，前者快而局部、后者慎而全局。

前沿方向上，2024--2025 年的工作集中在三处：零水平集附近的自适应采样以缓解高维稀疏、可验证的神经 BRT（训练后用 SMT/SOS 验证而非仅信损失）、以及环境变化时的在线 warm-start 更新。这些都在试图把“精确但算不动”与“可扩展但无保证”这对张力往中间推进。

% ════════════════════════════════════════════════════════════════
\section*{延伸阅读}

\begin{itemize}
  \item \textbf{综述与奠基}：Bansal、Chen、Herbert、Tomlin（2017，CDC）是最全面的入门综述\cite{bansal2017hamilton}；Mitchell、Bayen、Tomlin（2005，\emph{IEEE TAC}）是把可达集与时变 HJ 严格挂钩的工程奠基\cite{mitchell2005time}；Isaacs《微分博弈》是零和博弈源头\cite{isaacs1965differential}；Crandall--Lions（1983）黏性解\cite{crandall1983viscosity}、Osher--Sethian（1988）水平集\cite{osher1988fronts} 是两块数学地基。
  \item \textbf{CBVF 与统一}：Choi、Lee、Sreenath、Tomlin、Herbert（2021，CDC）是 CBVF 原文\cite{choi2021cbvf}；其 CBF 侧背景见 \cref{ch:clf-cbf}、HJB 侧见 \cref{ch:dp-hjb}。
  \item \textbf{高维与神经}：Bansal--Tomlin（DeepReach，ICRA 2021）\cite{bansal2021deepreach}；SIREN 架构见 Sitzmann 等（NeurIPS 2020）\cite{sitzmann2020siren}；PINN/自动微分基础见 \cref{ch:rl-basic}。
  \item \textbf{解耦框架与应用}：FaSTrack 见 Herbert 等（2017，CDC）\cite{herbert2017fastrack}；最小干预安全框架见 Fisac 等（2019，\emph{IEEE TAC}）\cite{fisac2019general}；多智能体分解见 Chen 等（\emph{IEEE TAC} 2018）\cite{chen2018decomposition}；reach--avoid 时变推广见 Fisac 等（2015，HSCC）\cite{fisac2015reach}。规划侧的 FaSTrack/RTD/Funnel 用途与微分博弈奠基分别见 \cref{sec:pr-reach,ch:plan-diffgame}。
\end{itemize}
